% Issue 12 draft, 24 November 2026. All four Discussion paragraphs approved.
% Issue 9.1: Zipf explanation approved after wording clarification.
% Includes the approved, explicitly labelled invented sentiment illustration.
% Total Discussion allowance: 0.55 pages; compiled fit remains for Issues 13--14.
\section*{Discussion}

Older-player transfers had a lower observed Premier League joining rate, with limited evidence against the origin-club/season null. In separate Guardian samples, older signings received more coverage, while tone changes showed no clear age difference. Greater coverage describes more attention, without establishing more negative treatment. These separate samples do not establish that press tone and recruitment decisions influence one another.

The lower observed entry rate does not establish that age itself restricted entry. Differences in player characteristics or preferences could contribute, but these explanations were not tested. Because the data record completed moves rather than offers and rejections, they cannot identify recruitment motives.

The sentiment comparison also depends on which signings had sufficient qualifying text. The filters retained 67 of 88 older signings and 753 of 1,244 younger signings, limiting generalisation to all selected signings. Player prominence and the mix of performances, topics and quoted voices could affect the measured tone, but their contributions were not tested. Some rare words may appear only once because our sample contains a limited amount of text. How we split expressions into tokens, and the repeated use of club names and injury updates, may also affect the frequency pattern. We did not test whether these factors explain why the vocabulary in our text analysis diverges from the ideal Zipf law in this log--log comparison of frequency rank and raw token frequency. RoBERTa scores the whole sentence without being told which player to evaluate. Its score may therefore reflect praise or criticism of another person mentioned in the same sentence \cite{hamborg2021news}. For an \textit{invented illustration}, consider ``Smith watched as Jones delivered an outstanding performance.'' This sentence names Smith but praises Jones, so a positive sentence score could enter Smith's coverage average without expressing praise for Smith. Validation by one annotator, including only 14 older-player sentences, leaves age-specific validity uncertain. Smaller differences in tone change remain possible; the inconclusive comparison does not establish equal changes.

Some players appear more than once, and the same article can contribute to several signings. The age-label shuffles may not fully account for these links between observations. This limits how confidently we can interpret the p-values \cite{hemerik2018exact,winkler2015multilevel}. The text findings describe selected English-language Guardian coverage. Monsees's comparison of German and Swedish newspapers motivates caution about generalising across editorial, linguistic and cultural contexts \cite{monsees2024talent}. Missing original article text and score caches also limit independent reconstruction.
